\chapter{Kidney Disease}

\section*{Patient focus}
Kidney imaging helps doctors find blockages, infections, cysts, tumours, and causes of reduced kidney function. Ultrasound is usually the first test because it is safe, painless, and does not use radiation. It can quickly tell if there is a blockage in the urine drainage system, measure kidney size, and check for cysts or masses. CT scans provide more detailed pictures when needed, and MRI can help characterise complex masses without radiation.


\label{ch:ch65_kidney_disease}

\section{Imaging}
Imaging in kidney disease assesses obstruction, infection, stones, cystic disease, masses, vascular disease, and treatment complications. Ultrasound, noncontrast CT, contrast CT, MRI, nuclear renography, and biopsy guidance have different roles.

\section{When to Image CKD}
\begin{itemize}
  \item \textbf{Acute kidney injury (AKI)}: ultrasound to exclude obstruction (obstructive uropathy); CT if stones suspected (non-contrast CT KUB)
  \item \textbf{Chronic kidney disease (CKD)}: imaging is not routinely required for diagnosis; focus on finding reversible causes
  \item \textbf{Indications for imaging in CKD}:
    \begin{enumerate}
      \item Suspected obstruction
      \item Recurrent infections
      \item Haematuria (to exclude mass)
      \item Asymmetric kidney size
      \item Rapid decline in function
      \item Planning for renal biopsy or transplantation
    \end{enumerate}
  \item \textbf{Kidney size}: normal 9--12 cm; small kidneys ($<$9 cm) suggest chronic disease; asymmetric size suggests renal artery stenosis or chronic pyelonephritis
\end{itemize}

\begin{pearlbox}[Kidney size matters]
Bilateral small, echogenic kidneys with thinned cortex are the classic appearance of end-stage CKD. A single small kidney with a normal contralateral kidney suggests renal artery stenosis on the small side---check with Doppler US or CTA.
\end{pearlbox}

\section{Polycystic Kidney Disease}

\subsection{Autosomal Dominant PKD (ADPKD)}
\begin{itemize}
  \item \textbf{US criteria for diagnosis} (age-dependent):
    \begin{enumerate}
      \item Age 15--39: $\geq$3 cysts (unilateral or bilateral)
      \item Age 40--59: $\geq$2 cysts in each kidney
      \item Age $\geq$60: $\geq$4 cysts in each kidney
    \end{enumerate}
  \item \textbf{CT/MRI}: shows innumerable cysts, kidney enlargement, may show hemorrhagic cysts (high attenuation on CT, high T1 signal on MRI)
  \item \textbf{Complications}: cyst infection, hemorrhage, stone disease, liver cysts, intracranial aneurysm
  \item \textbf{Prognosis}: total kidney volume (TKV) on imaging is a biomarker of disease progression---used in clinical trials (e.g.\ tolvaptan)
\end{itemize}

\subsection{Autosomal Recessive PKD (ARPKD)}
\begin{itemize}
  \item \textbf{Prenatal US}: enlarged, echogenic kidneys with loss of corticomedullary differentiation
  \item \textbf{Postnatal}: enlarged kidneys with multiple small cysts at corticomedullary junction, associated hepatic fibrosis
\end{itemize}

\section{Renal Artery Stenosis}
\begin{itemize}
  \item \textbf{Doppler US}: peak systolic velocity $>$180--200 cm/s at stenosis, renal aortic ratio $>$3.5; post-stenotic turbulence; small kidney with reduced resistive index
  \item \textbf{CTA}: spatial resolution allows detection of fibromuscular dysplasia (FMD) and atherosclerotic stenosis
  \item \textbf{MRA}: contrast-enhanced MRA for stenosis; non-contrast MRA (time-of-flight, phase contrast) if renal function precludes gadolinium
  \item \textbf{Gold standard}: conventional angiography---reserved for equivocal non-invasive results or when intervention is planned
\end{itemize}

\begin{table}[h]
\centering
\caption{Causes of renal artery stenosis}
\begin{tabular}{|l|l|l|}
\hline
\textbf{Cause} & \textbf{Age group} & \textbf{Imaging pattern} \\
\hline
Arteritis (Takayasu) & Young females & Long-segment, concentric \\
Atherosclerosis & $>$50 years & Proximal ostial stenosis, plaque \\
Fibromuscular dysplasia & Young females & Mid-distal, string-of-beads \\
Post-transplant & Any & Anastomotic stenosis \\
\hline
\end{tabular}
\end{table}

\section{Diabetic Nephropathy}
\begin{itemize}
  \item \textbf{Early disease}: US may be normal or show increased echogenicity
  \item \textbf{Progressive CKD}: bilateral small, echogenic kidneys with thinned cortex
  \item \textbf{Imaging role}: mainly to exclude other causes of renal decline (obstruction, stones, mass); not for staging diabetic nephropathy itself
  \item \textbf{Papillary necrosis}: more common in diabetics; CT shows cavities within renal papillae, contrast in collecting system
\end{itemize}

\section{Renal Transplant Imaging}

\subsection{Post-Transplant Ultrasound Protocol}
\begin{itemize}
  \item \textbf{Baseline (0--48 hours)}: confirm vascular patency (renal artery and vein Doppler), exclude perinephric collections, assess parenchymal echogenicity
  \item \textbf{Routine surveillance}: varies by institution; typically at 1, 3, 6, and 12 months
  \item \textbf{Indications for urgent US}: rising creatinine, reduced urine output, fever, graft tenderness
\end{itemize}

\subsection{Complications by Time Period}

\begin{table}[h]
\centering
\caption{Transplant complications by time period}
\begin{tabular}{|l|l|}
\hline
\textbf{Time period} & \textbf{Common complications} \\
\hline
1 week--3 months & Acute rejection, lymphocele, ureteric obstruction \\
3--12 months & Chronic rejection, BK virus nephropathy, CMV \\
$<$1 week & Vascular thrombosis, urine leak, haematoma \\
$>$1 year & Chronic allograft nephropathy, recurrent disease \\
\hline
\end{tabular}
\end{table}

\subsection{Key US Findings}
\begin{itemize}
  \item \textbf{Acute tubular necrosis}: oedematous kidney with increased resistive indices (RI $>$0.80)
  \item \textbf{Rejection}: increased RI, loss of corticomedullary differentiation, graft enlargement, may have perinephric fluid
  \item \textbf{Renal artery thrombosis}: absent arterial Doppler signal---surgical emergency
  \item \textbf{Renal vein thrombosis}: absent venous flow, graft swelling, reversed arterial diastolic flow
  \item \textbf{Ureteric obstruction}: hydronephrosis; may need antegrade nephrostomy
  \item \textbf{Lymphocele}: anechoic perinephric collection, may cause ureteric compression
\end{itemize}

\begin{pearlbox}[Resistive index in transplant]
RI $>$0.80 suggests raised intrarenal resistance---causes include acute rejection, acute tubular necrosis, nephrotoxic drugs, and urinary obstruction. RI alone cannot differentiate these causes; clinical correlation and often biopsy are needed.
\end{pearlbox}

\section{US Elastography for Renal Fibrosis}
\begin{itemize}
  \item \textbf{Shear wave elastography}: measures renal cortical stiffness in kPa or m/s; correlates with histological fibrosis
  \item \textbf{Applications}: staging CKD fibrosis, monitoring transplant fibrosis, reducing need for biopsy
  \item \textbf{Limitations}: technically challenging, operator-dependent, affected by hydronephrosis, acute inflammation, and recent biopsy
\end{itemize}

\section{Interpretation}
Describe kidney size, cortical thickness, echogenicity, collecting systems, perfusion, renal masses, cyst complexity, ureters, bladder, and perinephric tissue. Correlate with kidney function and urine findings.

\section{Safety}
Use contrast only after a patient-specific benefit-risk assessment; severe renal impairment does not automatically prohibit every contrast examination. Obstructed infection, hemorrhage, infarction, or acute vascular complication requires urgent escalation.

\begin{pearlbox}[Contrast in CKD]
An eGFR below 30 mL/min/1.73 m$^2$ or acute kidney injury calls for an individualized benefit-risk assessment; it is not an automatic prohibition. For intravenous iodinated contrast, consider isotonic IV saline prophylaxis when appropriate and not contraindicated. Use a contrast dose adequate for a diagnostic examination, and distinguish contrast-associated AKI from injury actually caused by contrast. The need for hydration depends on the patient's clinical status.
\end{pearlbox}

\clearpage
